\section{How to read this report}
\label{sec:glossary}

This report develops a mechanised semantics for a finite probabilistic
game-description language.  The intended reader has a working knowledge
of finite extensive-form games (players, information sets, mixed and
behavioural strategies, Nash equilibrium) but is not assumed to know
the modern factored-observation formalism, the Lean proof assistant,
or the specific corner of programming-language semantics that
underlies the construction.  We make three concessions to that mixed
audience.

\paragraph{Math first, Lean on the side.}
The body of the report uses ordinary mathematical notation.  Lean
identifiers appear only when their type or naming choice is itself the
subject of a design discussion; otherwise they are deferred to the
companion definitional supplement, which is the authoritative
cross-reference for every mathematical object introduced here.  Each
section ends with a brief \emph{Lean correspondence} note pointing to
the relevant supplement subsections.

\paragraph{Glossary.}
The following terms are used pervasively.  Readers from one
discipline can usually skip the entries that match their background.

\begin{description}[itemsep=0pt,leftmargin=1.2em,style=nextline]
\item[Finite-support distribution / probability mass function (PMF).]
  A function $\mu : \alpha \to [0,1]$ with finite support and
  $\sum_x \mu(x) = 1$.  We write $\PMF{\alpha}$ for the set of such
  functions.  $\PMF{\cdot}$ is a monad: the bind operator
  $\mu \bind f = \sum_x \mu(x)\,f(x)$ chains a stochastic step with a
  continuation indexed by its result, and the unit $\delta_x$
  (\emph{Dirac}) is the point mass at $x$.
\item[Finite-weight function.]
  An auxiliary unnormalised version, $\FWeight{\alpha} = \alpha \to_0
  \NNQ$, used internally during construction; converting to a PMF
  requires a side proof of total mass $1$.
\item[Kuhn's theorem (finite realisation form).]
  In a finite extensive-form game with sufficient observability,
  every mixed strategy over pure contingent plans induces the same
  outcome distribution as some behavioural strategy that randomises
  locally at each information set~\citep{Kuhn1953}.  The classical
  statement requires \emph{perfect recall}.  In the current Lean
  development the bounded-FOSG theorem is unconditional on the FOSG
  carriers; the native Vegas kernel-game statement is exposed with an
  explicit native/FOSG strategy-equivalence bridge obligation.
\item[Factored-observation stochastic game (FOSG).]
  A modern state-based formalism for finite extensive-form games with
  imperfect information~\citep{Kovarik2022}: each round, an active set
  of players picks actions from local alphabets, the state transitions
  stochastically, and per-player observations are emitted.  FOSGs
  separate cleanly from extensive-form trees by indexing strategies on
  per-player observation sequences rather than on global histories.
\item[Extensive-form game (EFG).]
  The classical tree presentation, with chance and decision nodes,
  per-player information sets, and outcomes at the leaves.
\item[Kernel game.]
  A minimal strategic-form carrier sufficient to define expected utility
  and Nash equilibrium: a per-player strategy type, an outcome type, a
  stochastic kernel from joint strategies to outcomes, and a utility
  $u : \Outcome \to \Players \to \mathbb{R}$.  Distinct from a
  normal-form game in that strategies need not be hand-written
  actions: they may be functions from observations to actions, and
  outcomes may carry more than a payoff vector.
\item[Static-single-assignment (SSA).]
  A standard programming-language convention~\citep{cytron1991efficiently}
  in which every variable is assigned at exactly one program point.  A
  program in SSA form has a unique definition site for each name; we
  exploit this to make every event node carry a unique write target.
\item[Event structure / dependency graph.]
  In concurrency theory~\citep{NielsenPlotkinWinskel1981,Winskel1986EventStructures}
  an event structure abstracts a process by its events and a causal
  dependency relation; independent events are unordered and may be
  permuted without affecting the resulting state.  Our event graphs
  are a typed, finite, asymmetric variant adapted to the protocol
  setting.
\item[Refinement / forward simulation.]
  A binary relation between an implementation and a specification
  asserting that the implementation cannot be distinguished from the
  specification by any externally observable behaviour.  Conventionally
  proved by exhibiting a state projection and matching transitions
  step-by-step~\citep{LynchVaandrager1995,AbadiLamport1991}.
\item[Visibility / view.]
  The information-set structure of a multi-party protocol, presented
  here as a binding-level discipline: every variable is either public
  or hidden to a single owner, and every program point exposes a
  per-player projected context.
\end{description}

\paragraph{Lean correspondence.}
For a top-down map from the bird's-eye picture to file paths, see
supplement~\S1.  Cross-references in this report use the supplement's
\TT{Vegas.X} naming.  Every theorem and definition stated here is
mechanised; the supplement records the exact Lean type and field
list.
